% !TeX root = ../papers/nc-p3.tex
\section{二つのねじれた階数一直線と方向の二次曲面}
\label{sec:rank-one-skew-pair-exclusion}
\begin{theorem}\label{thm:rank-one-skew-pair-exclusion}
被約な四辺有限の標準有限Frobenius窓の一般可逆な全階数反対根では、
階数一曲線の台は空集合、一本の直線、整平面二次曲線、交わる二直線のいずれかである。
\end{theorem}
\begin{proof}[方向と繊維の式]
二本の素な直線があるとし、それらを相補な二次元空間の射影直線に取る。
出力の対応する二直線も素で、各直線では平方の一次座標を入力と合わせる。
両直線外の入力を $(s u,t v)$ と書く。$u,v$ はそれぞれ二次元列、方向 $([u],[v])$ は
Segre二次曲面 $S$ 上の点である。各方向の横断直線上の直線外印は二つ以下である。
実際四次小行列式は両端で零となり、恒等零なら前定理の総次数上限に反する。
二直線上に各三印までしかないから、直線外印は十四以上必要である。

二直線を消す二次式は混合四成分に尽き、平方座標同型はこの四空間を同型に移す。
従って可逆な $T$ と二つの対称二行の双一次行列 $D,E$ により、対角ブロックは
$s^2uu^t+stD(u,v)$ と $t^2vv^t+stE(u,v)$、非対角ブロックは $st T(u\otimes v)$ である。
後者の方向が階数一であるところで $T(u\otimes v)=a\otimes b$ と書く。
$\rho=s/t$ とすると、直線外で階数一となる必要十分条件は、非零な $c$ があって
\[
        \rho uu^t+D=c aa^t,\qquad \rho^{-1}vv^t+E=c^{-1}bb^t
\]
となることである。以下で同じ文字の二乗は列の対称平方を表す。
\end{proof}
\begin{proof}[曲線上の点数]
両Segre上の方向を持つ既約曲線の正規化で、四つの射影方向の線束を $P,Q,A,B$、
その次数を $p,q,r,k$ とする。同じ超平面束 $H=P+Q=A+B$ の次数を $d$ とする（加法記法）。
この同型で上のテンソル分解が成り立つように取る。
$\delta_1=u\wedge a,\ \delta_2=v\wedge b$ の次数は $e_1=p+r,e_2=q+k$、和は $2d$ である。
二つとも恒等零でなければ、この曲線上の直線外印は $3d$ 以下である。以下、特異な方向では
一つの逆像を選んで数えればよい。$Jz=(-z_2,z_1)^t$ と書くと全ての採用方向で
\[
        f=(Ju)^tD(Ja)=0,\qquad g=(Jv)^tE(Jb)=0 .
\]
その次数は $d+e_1,d+e_2$。いずれかが非零なら方向数は大きくても $d+\max(e_1,e_2)$。
二つの異なる $\rho$ が解なら二式の差から $\delta_1=\delta_2=0$ だから、
重ねて数える分は $\min(e_1,e_2)$ 以下である。

$f,g$ が恒等零なら有理的に $D=c_1u^2+c_2a^2$, $E=c_3v^2+c_4b^2$ と一意に書ける。
係数の極はそれぞれ $\delta_1,\delta_2$ の零因子までである。
実際二つの平方列の外積は元の二列の行列式と、至る所で非零な二列の対称積を定数同型で移したものとの積である。
局所離散付値環で単元となる成分を使って係数を解けばよい。
さらに一致方向の採用点では元の行列もその共通平方に比例するので、分母を払った二係数はともに零になる。
一般の方向に解がある条件は
\[
                         c_1 c_3=1,\qquad c_2 c_4=1 .
\]
各積は関数である。二式とも恒等式なら $\rho=-c_1$ で第三の階数一曲線が生じる。
従って少なくとも一つは恒等式でなく、その分子は次数 $e_1+e_2=2d$ の非零切断である。
一致方向の採用点も上の分子消滅により数えられる。方向数は $2d$ 以下、
重複分込みで $3d$ 以下となる。
\end{proof}
\begin{proof}[一方の方向が曲線全体で一致する場合]
例えば $a=u$ とし、$A=P,B=Q$ と同定する。$v,b$ は一般には異なるとする。
$D Ju$ が恒等零でなければ方向数は $d+p$ 以下、重複分は $2q$ 以下である。
恒等零なら $D=\alpha u^2$ で $\alpha$ は $Q-P$ の正則切断である。
$g\ne0$ なら方向は $d+2q$ 以下である。$g=0$ が恒等的なら上の係数 $c_3,c_4$ を用い、
$v\ne b$（射影方向）の採用点では $\rho^{-1}=-c_3$ と
$(\rho+\alpha)c_4=1$ を得る。よって
\[
                       (\alpha c_3-1)c_4-c_3=0 .
\]
左辺は $P-Q$ の有理切断で極は $2\operatorname{div}_0(\delta_2)$ まで、
分子の次数は $p+3q=d+2q$。一致方向の採用点でも分子は零である。
恒等零で $c_3,c_4$ が一般に非零なら第三の曲線ができる。一方が恒等零なら両方が零で、
採用方向は一致する $2q$ 点以下に限られる。従って $D=\alpha u^2$ の場合は
重複分を足しても $d+4q$ 以下である。入出力二ブロックを交換しても同じ議論が使える。

方向の直線成分では常に三印以下となる。一致が恒等的にない場合は前段を使う。
入出力で異なる側の因子が固定なら両方の一致とも恒等的になり得ない。
同じ側の因子が固定なら、該当する印は対応する一対の平面 $L,L'$ に属し、
そこには他方の座標軸直線とその像も含まれる。
二つの平面の平方の包は平方座標同型で一致できない。さもなければ $L'$ の零化列が $L$ で消え、
定理\ref{thm:opposite-column-divisor}に反する。従って該当印は $L$ の非零二次式上にあり、
軸直線を除くと残余直線上の三つまでである。

方向の整二次成分では $p=q=r=k=1$ で、片側一致でも上界は六である。
両側一致なら、$D Ju,E Jv$ のいずれかが恒等零でない場合はその次数三によって三方向、各二印までである。
両者とも恒等零なら $D=\alpha u^2,E=\beta v^2$。$\alpha,\beta$ は正規化で互いに逆次数の
次数零の束の切断なので各々恒等零か至る所で非零である。共通テンソルを $u\otimes v$ に合わせると
残る式は $(\rho+\alpha)(\rho^{-1}+\beta)=1$（$a=u,b=v$ の規格化による定数倍を右辺に許す）。
一般点で非零解があれば第三の曲線ができ、なければ次数零の束を自明化した同じ定数方程式が
どの点でも解を持たない。
\end{proof}
\begin{proof}[異なる二次曲面の交線]
$T(S)\ne S$ なら方向は $(2,2)$ 因子 $S\cap T^{-1}S$ に入る。整四次なら
片側一致が恒等的に成り立つと全印で対応する非零混合交代式が消える（両軸上でも消える）。
従って上の $3d$ 評価で十二印以下となる。
全成分が直線と二次曲線なら、重複成分があっても被約次数の和は四以下なので同じ結論である。

整三次成分 $K$ があれば他は直線 $R$ 一本で、$K$ は滑らかな $(1,2)$ または $(2,1)$ 曲線である
（一方の次数一の射影でグラフとなる）。一致が恒等的になければ $K$ は九印以下である。
両側一致なら $T$ は $K$ の全点を射影的に固定し、$K$ は非平面であるから $T$ はスカラーとなる。
片側のみの一致を $u=a$ とする。$p=2,q=1$ なら先の二上界 $d+p+2q,d+4q$ はともに七である。
$p=1,q=2$ で $D Ju$ が恒等零でなければ八以下である。
残る $D=\alpha u^2$ の場合も $K$ には十一印までである。このとき $R$ 上の $K$ 外の印は一つ以下であることを示す。

$R$ 上は $v,b$ が固定され、$a$ は $u$ の射影同型 $M$ で動く。$M$ がスカラーなら全印で $a=u$ となり矛盾。
$\delta=u\wedge Mu$ の零因子は $Z=K\cap R$ である。実際 $K$ 上での恒等的一致により
$\delta|_Z=0$ であり、両方の長さは二である。$D|_R$ は $u$ の一次行列であり、
$Z$ 上では $D$ は $u^2$ に比例する（非被約な交叉でも $K$ 上の等式を制限する）。
従って次数三の $f|_R$ は $Z$ で零となる。これが非零なら $Z$ 外で一方向までで、
そこは $u\ne a$ なので一印までである。恒等零なら $D=c_1 u^2+c_2(Mu)^2$ と分解した二係数は
$Z$ を含め正則である。実際 $D\bmod u^2$ は局所的に $Z$ の式で割れ、
$(Mu)^2\bmod u^2$ はその式と原始列の積なので、$c_2$ は極を持たず $c_1$ も同様である。
両者の束は $\mathcal O_R(-1)$ だからともに零になり、$Z$ 外に解はない。
これで三次を含む場合も十二印以下となり、十四に届かない。
\end{proof}
\begin{proof}[Segre二次曲面の自己同型]
二族の母線を交換する場合は例えば $a$ が $v$ の固定射影変更であり、
両軸を含む全印で非零混合交代式が消えて不可能である。二族を保つ場合は
$T=M\otimes L$ とし、$a=Mu,b=Lv$ と取れる。両行列とも非スカラーである。
片方がスカラーなら同じ側の混合交代式が全印で消える。
方向曲面全体で $f,g$ はそれぞれ二次数 $(3,1),(1,3)$ の非零式である。
実際 $f$ が恒等零なら上の分解の $c_1$ に対し $F=(u\wedge Mu)c_1$ は二次数 $(1,1)$ の正則切断となる。
平方列の外積による分母評価はここでも各素因子で成立し、滑らかな曲面なので大域的に正則である。
実際の入力二列を $u,v$ とすれば $u\wedge Mu+F(u,v)$ は非零な斉次二次式である。
直線外印で $u\ne Mu$ なら $c_1=-1$、一致方向でも $D$ が共通平方に比例するため $F=0$。
また $u=0$ 上でもこの式は零であり、二十印のうち $v=0$ の三つまででしか非零になれない。
これは非零二次式の評価の支持下界六（定理\ref{thm:reducible-octic-factors}の二次式の跡評価）に反する。
$g$ も同様である。

$f,g$ の共通因子がなければ交点長は十で、重複計数できる方向は両行列の固有方向の対の高々四つ。
十四印なら全四対に二印ずつあり、両軸に各三印を要する。
各第一固有方向を固定した平面は第二軸の三印と既に四つの直線外印を持ち、
他の印を持てない。第二固有方向でも同じである。各第一固有方向 $u_i$ で
二つの独立な $v_j$ に対し $D(u_i,v_j)$ が $u_i^2$ に比例するから、双一次性により全ての $v$ で比例する。
従って $f$ は二つの第一固有方向の繊維を含み、$g$ も二つの第二固有方向の繊維を含む。
これらの外で残る二式は $(1,1)$ で固有に交わるため二方向以下だが、そこには六印が必要である。

最大共通因子が $(1,0)$ か $(0,1)$ ならその上は三印以下、残余交点長は六で、重複四つを加えても十三以下。
共通因子が $(1,1)$ ならその上は六印以下、外の交点長は四なので十四に達するには
全四固有対が共通因子の外にあって各二印を持つ必要がある。この場合、共通部分のどの成分でも
一般点で両一致はともに成立しない（片方が至る所一致する成分は固有方向の繊維で、固有対を含む）。
各成分で $f=g=0$ が恒等的であるから、その上の採用方向数は既に示した $2d$ 以下。
共通部分の総次数は二で、そこには重複もなく、四印までしかない。矛盾である。
最大共通因子の二次数はこれらで尽き、全ての場合に素な二直線を排除した。
\end{proof}
